\begin{proposition}\label{prop:octic-cyclic-realizations}
The groups
\[
 C_{33},\quad C_{34},\quad C_{35},\quad C_{39},\quad C_{40},\quad C_{42}
\]
belong to $\Phi(8)$.
\end{proposition}
\begin{proof}
We verify all ten degree-eight entries in van Hoeij's
\texttt{LowDegreePlaces} data. Their levels, in source order, are
\[
 33,34,34,34,35,39,40,42,42,42.
\]
The exact polynomials and coordinates are preserved in
\texttt{inputs/vanhoeij\_octic\_input.json} and in the corresponding
Magma input file, with the original source lines and their hash.
For each entry let $(x,y)$ be its listed coordinate pair and put
\[
 r=\frac{x^2y-xy+y-1}{x^2y-x},\qquad
 s=\frac{xy-y+1}{xy},\qquad b=rs(r-1),\qquad c=s(r-1).
\]
Over its octic field $K$, define
\[
 E:\quad Y^2+(1-c)XY-bY=X^3-bX^2,\qquad P=(0,0).
\]
Every field polynomial is checked irreducible of degree $8$.
For the level-$33$ entry, originally given by a quartic equation
in $x$ and a quadratic equation in $y$, the primitive generator
$a=x+y$ has polynomial
\[
 a^8-15a^7+95a^6-331a^5+741a^4-1100a^3+1009a^2-688a+359.
\]
Exact elimination recovers both coordinates as polynomials in $a$,
and the two original equations are checked again over $K$.

For the entry at level $N$, exact arithmetic verifies
$NP=0$ and $(N/q)P\ne0$ for every prime divisor $q$ of $N$.
The following good reductions certify the full torsion. A triple
$(p,f,h)$ records the rational residue characteristic, residue
degree and finite elliptic-curve order.
\[
\begin{array}{c|r|l|r}
\text{entry}&N& (p,f,h)&\gcd h\\\hline
1&33&(7,8,5760579),(13,8,815675619)&33\\
2&34&(5,2,34)&34\\
3&34&(7,4,2448),(11,2,102),(13,4,28900)&34\\
4&34&(7,4,2482),(13,4,28424)&34\\
5&35&(7,8,5763520),(11,2,105)&35\\
6&39&(7,2,39)&39\\
7&40&(3,4,80),(31,3,29560)&40\\
8&42&(41,1,42)&42\\
9&42&(7,8,5765886),(13,2,168)&42\\
10&42&(7,8,5765886),(13,2,168)&42
\end{array}
\]
All displayed primes are odd. The field polynomial is monic and
integral at the given prime, with unit discriminant; hence the
prime is unramified in $K$. All coordinate denominators are units
there, and the reduced Tate equation has nonzero discriminant.
Reduction is therefore injective on the entire torsion subgroup.
Its order divides the displayed gcd, which the point $P$ attains.
Thus every entry has full torsion $C_N$.
The complete source and log are
\texttt{verify\_octic\_examples.m} and
\texttt{verify\_octic\_examples.log}.
\end{proof}

\begin{corollary}\label{cor:octic-known-isogeny-bounds}
For each of these ten curves, every elliptic curve isogenous over
the same octic field has torsion order dividing the same $N$.
In particular, these isogeny classes cannot produce an additional
torsion group outside the six groups above and the published
$\Phi^\infty(8)$ list.
\end{corollary}
\begin{proof}
Good reduction is invariant under isogeny. Isogenous curves over
a finite field have the same number of points, so the same gcd
bounds apply. The numbers $33,34,35,39,42$ are squarefree, so a
subgroup of $E(\overline K)_{\mathrm{tors}}$ with order dividing
one of these numbers is cyclic. For $N=40$, the only noncyclic
possibilities have the form $C_2\times C_{2d}$ with $d\mid10$;
all belong to $\Phi^\infty(8)$. Every remaining cyclic divisor
is either already in the published infinite list or is one of
the six groups in the proposition.
\end{proof}

% These are realization certificates. The probabilistic source search is
% not used as evidence of completeness of the octic classification.
